\documentclass[12pt,a4paper]{article}
\usepackage[a4paper, top=2.5cm, bottom=2.5cm, left=2.8cm, right=2.8cm]{geometry}
\usepackage{times}
\usepackage[utf8]{inputenc}
\usepackage[T1]{fontenc}
\usepackage{graphicx}
\usepackage{amsmath}
\usepackage{booktabs}
\usepackage{caption}
\usepackage{hyperref}
\usepackage{xcolor}
\usepackage{titlesec}
\usepackage{parskip}
\usepackage{enumitem}

\hypersetup{colorlinks=true, linkcolor=blue, citecolor=blue, urlcolor=blue}
\titleformat{\section}{\large\bfseries}{}{0em}{}[\titlerule]
\titlespacing{\section}{0pt}{14pt}{6pt}
\setlength{\parindent}{0pt}
\setlength{\parskip}{6pt}

\newcommand{\placeholderimg}[2]{%
  \fbox{\begin{minipage}{#1}\centering\vspace{0.8cm}%
    {\color{gray}\small\texttt{[INSERT: #2]}}\vspace{0.8cm}%
  \end{minipage}}%
}

\begin{document}

\begin{center}
  {\LARGE\bfseries Lung Segmentation on the COVID-QU-Ex Dataset\\[4pt]
  Using a U-Net Deep Learning Model\par}
  \vspace{0.5cm}
  {\large Nguyen Xuan Quan \quad Student ID: 23BI14370\par}
\end{center}
\vspace{0.4cm}\hrule\vspace{0.4cm}

\section{Introduction}

The COVID-QU-Ex dataset, compiled by researchers at Qatar University, is one of
the largest publicly available chest X-ray (CXR) repositories for COVID-19
research. It contains \textbf{33,920} CXR images across three diagnostic
categories: \textbf{COVID-19} (11,956 images), \textbf{Non-COVID infections}
(viral or bacterial pneumonia, 11,263 images), and \textbf{Normal} (10,701
images). Every image is paired with a ground-truth \emph{lung segmentation mask},
making this the largest lung mask dataset ever created.

Accurate lung segmentation is a prerequisite for downstream classification tasks:
it restricts the model's attention to the pulmonary region and discards irrelevant
background anatomy (ribs, soft tissue, medical devices). This project applies a
U-Net convolutional neural network to perform pixel-level binary segmentation of
the lung region from CXR images.

\section{Exploratory Data Analysis (EDA)}

\subsection*{2.1 Directory Audit and File Count}

The first two code cells define a \texttt{count\_files\_in\_directories} function
that walks the entire dataset directory tree and counts the \texttt{.png} files
present in each sub-folder (Train / Val / Test, each containing three class
sub-directories). This is a critical quality-control step: a mismatch between
image count and mask count would silently corrupt training by pairing the wrong
mask with the wrong X-ray. The audit confirmed a one-to-one correspondence across
all splits and categories.

\subsection*{2.2 Visual Inspection of Image--Mask Pairs}

Cell 3 renders a sample lung mask from the Normal class. Cells 7 and 8 then
display a training X-ray and its corresponding binary mask side by side. These
visual alignment checks confirm that lung boundaries in the image coincide exactly
with the white (foreground) regions in the mask, validating annotation quality
before any training begins.

\begin{figure}
    \centering
    \includegraphics[width=0.5\linewidth]{output.png}
    \caption{X-ray}
    \label{fig:placeholder}
\end{figure}
\begin{figure}
    \centering
    \includegraphics[width=0.5\linewidth]{output2.png}
    \caption{Mak}
    \label{fig:placeholder}
\end{figure}
\subsection*{2.3 Pixel-Value Distribution of Predicted Masks}

Cell 19 executes \texttt{np.unique(predict\_masks)} on the model output. The
result shows values concentrated firmly near 0 and near 1:

\begin{center}
\small\texttt{array([0.00539569, 0.00559962, \ldots, 0.9999989, 0.9999993],
dtype=float32)}
\end{center}

This bimodal distribution is the hallmark of a well-trained binary segmentation
model. The sigmoid output has learned to commit decisively to ``background''
($\approx 0$) or ``lung'' ($\approx 1$), rather than remaining uncertain around
0.5. This confirms the Dice loss and sigmoid activation are cooperating correctly.

\section{Data Preprocessing}

\subsection*{3.1 Image Loading and Resizing to 256$\times$256}

The \texttt{load\_images\_from\_folder} function (Cell 5) reads each \texttt{.png}
in grayscale mode and reshapes it to \texttt{(256, 256, 1)}. The U-Net requires
all inputs to share identical spatial dimensions. A resolution of 256$\times$256
balances image detail (clinically relevant features such as pleural boundaries
are preserved) with GPU memory efficiency (batch size of 16 fits comfortably).

\begin{center}
\begin{tabular}{lcc}
\toprule
\textbf{Split} & \textbf{Images loaded} & \textbf{Masks loaded} \\
\midrule
Training   & 3,000 & 3,000 \\
Validation & 200   & 200   \\
Test       & 50    & 50    \\
\bottomrule
\end{tabular}
\end{center}

\subsection*{3.2 Normalisation and Binarisation}

The \texttt{normalize\_and\_convert\_to\_binary} function (Cell 12) applies:

\textbf{(a) Image normalisation:}
\[
  x_{\text{norm}} = \frac{x - 255}{255} \quad \Rightarrow \quad
  x_{\text{norm}} \in [-1,\; 0]
\]
Normalisation prevents large raw pixel values (0--255) from generating
disproportionately large gradient magnitudes in the early convolutional layers,
which would destabilise or slow convergence. Batch Normalisation layers inside
the model subsequently re-centre the feature distributions during forward passes.

\textbf{(b) Mask binarisation:}
\[
  m_{\text{binary}} = \mathbb{1}[m > 127] \in \{0.0,\; 1.0\}
\]
Thresholding at 127 removes any JPEG compression artefacts near mask edges and
produces the clean binary targets required by the Dice loss function.

% ═══════════════════════════════════════════════════════════════════════════
% IMAGE NOTE 2 — chèn ở đây
% Lấy từ: Cell 14 (normalized X-ray) và Cell 15 (binary mask) trong Jupyter
% Cách lấy: chạy lại Cell 14 → right-click → Save as xray_norm.png
%            chạy lại Cell 15 → right-click → Save as mask_binary.png
% Thay \placeholderimg bằng:
%   \includegraphics[width=0.44\linewidth]{xray_norm.png}\hfill
%   \includegraphics[width=0.44\linewidth]{mask_binary.png}
% ═══════════════════════════════════════════════════════════════════════════
\begin{figure}
    \centering
    \includegraphics[width=0.5\linewidth]{output3.png}
    \caption{Normalized X-ray}
    \label{fig:placeholder}
\end{figure}

\begin{figure}
    \centering
    \includegraphics[width=0.5\linewidth]{output4.png}
    \caption{Ground-truth binary mask}
    \label{fig:placeholder}
\end{figure}
\section{Model Architecture}

\subsection*{4.1 U-Net Overview}

The U-Net (Ronneberger et al., 2015) is the standard architecture for biomedical
image segmentation. Its ``U'' shape consists of three components:

\begin{enumerate}[noitemsep]
  \item \textbf{Encoder:} four convolution + max-pooling blocks reduce the
    spatial size from 256$\times$256 down to 16$\times$16 while increasing
    channel depth (32 $\to$ 64 $\to$ 128 $\to$ 256).
  \item \textbf{Bottleneck:} a 512-filter convolution block at the lowest
    resolution encodes a compressed global representation.
  \item \textbf{Decoder:} four transposed-convolution blocks mirror the encoder,
    restoring the spatial size back to 256$\times$256 to produce the final mask.
\end{enumerate}

\textbf{Skip connections} concatenate encoder feature maps directly to the
corresponding decoder layers. This is the key innovation: fine spatial
details (exact lung boundary positions) lost during downsampling are passed
directly to the decoder, enabling sharp, well-localised segmentation boundaries.

\subsection*{4.2 Implementation Details}

Each \texttt{conv\_block} applies: Conv2D $\to$ BatchNorm $\to$ ReLU $\to$
Dropout(0.2) twice in sequence. Key design decisions:

\textbf{BatchNormalization} after each convolution accelerates convergence by
reducing internal covariate shift and acts as mild regularisation.

\textbf{Dropout(0.2)} drops 20\% of activations randomly during training,
preventing co-adaptation of neurons and reducing overfitting --- important
with only 3,000 training images.

\textbf{Sigmoid output} at the final \texttt{Conv2D(1, 1)} maps each pixel
to a probability in $[0, 1]$.

\textbf{Total parameters:} 7,771,297 (7,765,409 trainable; 5,888 non-trainable
from BatchNorm running statistics).

\subsection*{4.3 Loss Function and Metrics}

\textbf{Dice Loss} was chosen over binary cross-entropy because it directly
optimises the spatial overlap between prediction and ground truth, making it
inherently robust to class imbalance (background pixels far outnumber lung pixels):
\[
  \mathcal{L}_{\text{Dice}} = -\frac{2 \sum y_{\text{true}} \cdot y_{\text{pred}}
  + 1}{\sum y_{\text{true}} + \sum y_{\text{pred}} + 1}
\]

The \textbf{Jaccard Index} (Intersection-over-Union) was monitored alongside:
\[
  J = \frac{|A \cap B|}{|A \cup B|}
\]

\textbf{Adam} with learning rate $10^{-5}$ was used. The small learning rate
produces stable, incremental updates appropriate for a relatively small dataset.

\section{Model Training}

\subsection*{5.1 Configuration}

Training ran for 100 full epochs with batch size 16. Three callbacks were
registered: \texttt{ModelCheckpoint} (saves best weights by \texttt{val\_loss}),
\texttt{ReduceLROnPlateau} (halves LR after 10 stagnant epochs), and
\texttt{EarlyStopping} (patience 30). Neither the LR scheduler nor early stopping
was triggered across 100 epochs, confirming the learning rate of $10^{-5}$
remained well-suited throughout training.

\subsection*{5.2 Learning Curves}

\begin{figure}[h]
  \centering
  \includegraphics[width=\linewidth]{output5.png}
  \caption{Training and validation performance over 100 epochs (Cell 13).
  \textbf{Left:} Dice loss (negative Dice coefficient) decreases steadily 
  from $-0.43$ at epoch 1 to $-0.96$ at epoch 100 for both splits. 
  \textbf{Right:} Dice coefficient increases from $\approx 0.43$ to 
  $\approx 0.96$. The validation curve (orange) tracks the training curve 
  (blue) closely throughout, with minor fluctuations caused by the small 
  validation set size (200 images). The absence of divergence between the 
  two curves confirms the model generalises well without overfitting.}
  \label{fig:curves}
\end{figure}
The learning curves reveal several important observations. Both training and
validation Dice increase smoothly and consistently, a sign of a stable learning
rate and good model initialisation. The final training Dice (0.9634) and
validation Dice (0.9581) diverge by only 0.005 points, confirming that the
Dropout and BatchNorm regularisation effectively prevented overfitting despite
the small training set. Occasional fluctuations in the validation curve
(e.g., epochs 15 and 67) are expected given the small validation set of only
200 images, which produces higher per-epoch variance.

\section{Results and Evaluation}

\subsection*{6.1 Final Metrics at Epoch 100}

\begin{center}
\begin{tabular}{lcc}
\toprule
\textbf{Metric}      & \textbf{Training} & \textbf{Validation} \\
\midrule
Dice Coefficient      & 0.9634            & 0.9581              \\
Jaccard Index (IoU)   & 0.9849            & 0.9848              \\
Binary Accuracy       & 0.9929            & 0.9881              \\
\bottomrule
\end{tabular}
\end{center}

A Dice coefficient of \textbf{0.958} means that on average 95.8\% of the
predicted lung pixels perfectly overlap with the ground-truth annotation ---
considered excellent performance for medical image segmentation. The nearly
identical Jaccard value (0.985) is consistent with the Dice score via the
relation $J = D / (2 - D)$.

\subsection*{6.2 Qualitative Prediction}
\begin{figure}[h]
  \centering
  \includegraphics[width=0.44\linewidth]{output6.png}\hfill
  \includegraphics[width=0.44\linewidth]{output7.png}
  \caption{Qualitative segmentation result on a held-out test image (index 7).
  \textbf{Left:} Normalised chest X-ray (Cell 17) passed to the trained U-Net
  as input. \textbf{Right:} Corresponding lung mask predicted by the model
  (Cell 18). The predicted mask cleanly delineates both left and right lung
  lobes with sharp, well-defined boundaries, consistent with the high Dice
  coefficient of 0.958 achieved on the validation set.}
  \label{fig:prediction}
\end{figure}
The bimodal output probability distribution (Section 2.3, Cell 19) provides
additional quantitative validation: confident predictions near 0 and 1 indicate
the model has learned a sharp decision boundary rather than producing uncertain
mid-range probabilities.

\section{Discussion}

\textbf{Strengths.} The model achieves near state-of-the-art segmentation accuracy
using only 3,000 training samples and 100 epochs. The U-Net's skip connections
proved critical for preserving fine edge detail at the lung boundaries, while the
Dice loss handled the inherent class imbalance effectively. Training and validation
metrics diverged by less than 0.006 Dice points, demonstrating robust generalisation.

\textbf{Limitations.} The experiment was restricted to the Normal class due to
computational constraints. Lung morphology in COVID-19 and Pneumonia cases
differs substantially from normal lungs (ground-glass opacities, consolidations,
and pleural effusions can distort boundaries), so the high metrics reported here
may not transfer directly to pathological classes. Additionally, the normalisation
formula $(x - 255)/255$ maps values to $[-1, 0]$ rather than the more
conventional $[0, 1]$; while the model learns successfully regardless, correcting
this would improve interpretability and compatibility with pre-trained weights.

\textbf{Future Work.} Training on all three classes, applying data augmentation
(random flips, rotations, elastic deformations), and post-processing with
connected-component filtering to remove spurious small predicted regions would
all be natural extensions.

\section{Conclusion}

This project demonstrated that a U-Net with Dice loss, Batch Normalisation, and
Dropout achieves excellent lung segmentation on the COVID-QU-Ex Normal class.
Starting from random initialisation, the model reached a Dice coefficient of
\textbf{0.958} and a Jaccard index of \textbf{0.985} on the validation set after
100 epochs of training on 3,000 chest X-rays. The closely tracking learning
curves confirm stable generalisation without overfitting. These results establish
a strong foundation for extending the pipeline to pathological CXR classes and
integrating lung segmentation as a preprocessing stage for downstream COVID-19
classification.

\section*{References}

\begin{enumerate}[noitemsep]
  \item Ronneberger, O., Fischer, P., \& Brox, T. (2015). U-Net: Convolutional
    Networks for Biomedical Image Segmentation. \textit{MICCAI 2015}, LNCS 9351,
    pp.\ 234--241.
  \item Tahir, A. M., et al.\ (2022). COVID-QU-Ex Dataset. \textit{IEEE Dataport}.
    \url{https://doi.org/10.21227/25ex-p371}
  \item Dice, L. R. (1945). Measures of the Amount of Ecologic Association Between
    Species. \textit{Ecology}, 26(3), 297--302.
  \item Kingma, D. P., \& Ba, J. (2014). Adam: A Method for Stochastic
    Optimization. \textit{arXiv:1412.6980}.
\end{enumerate}

\end{document}